\section{Robustez e diagnóstico dos modelos}
\label{sec:robustez}

Três baterias de verificação fecham a fase de estimação.

\textbf{Sem o biênio pandêmico.} Todos os modelos foram reestimados
excluindo 2020 e 2021, os anos em que a casuística, os denominadores e
o cadastro de leitos sofreram os choques documentados na primeira
parte. A Figura~\ref{fig:robustezpandemia} compara os contrastes OSS vs
Direta nas duas amostras: as diferenças são mínimas (o tempo de
permanência passa de 10,2\% para 11,3\% de queda, a produção de 45,4\%
para 46,6\% de alta, e assim por diante). \textbf{Nenhuma conclusão
deste documento depende dos anos de pandemia.}

\begin{figure}[H]
\centering
\includegraphics[width=.96\textwidth]{fig_est_05_robustez_sem_pandemia.png}
\caption{Contraste OSS vs Direta estimado no painel completo e na
amostra sem 2020 e 2021, para os modelos em escala percentual.}
\label{fig:robustezpandemia}
\end{figure}

\comoler{cada linha traz duas estimativas do mesmo efeito: o círculo
azul escuro usa todos os anos e o quadrado verde exclui 2020 e 2021.
Quanto mais próximos os dois símbolos, mais robusto o resultado; aqui
eles praticamente se sobrepõem em todas as linhas.}

\textbf{Mortalidade ajustada.} O modelo de mortalidade foi refeito com
a versão que exclui óbitos fetais, maternos e neonatais: o efeito OSS
within vai de 0,230 para 0,233 negativos no logito e o efeito marginal
médio de 1,12 para 1,14 ponto percentual, confirmando que as duas
versões do indicador contam a mesma história, como a correlação
perfeita da primeira parte já anunciava.

\textbf{Resíduos quantílicos.} A primeira parte prometeu que a escolha
das famílias de distribuição seria validada nos modelos ajustados, e é
o que as Figuras~\ref{fig:qqprimeira} e~\ref{fig:qqsegunda} fazem. O
resíduo quantílico é um teste de qualidade: se o modelo descreve bem os
dados, esses resíduos se alinham sobre a diagonal do gráfico. O
alinhamento é muito bom na mortalidade (correlação de 0,984), na
fração de alta complexidade (0,991) e na produção (0,968), e bom no
custo (0,970) e no tempo de permanência (0,952), cujas caudas pesadas
já estavam diagnosticadas e são absorvidas pelas margens de erro
agrupadas. A exceção honesta é a \textbf{intensidade de ocupação da
UTI} (0,814): a cauda de UTIs com uso quase nulo escapa do modelo em
logaritmo, e esse resultado específico deve ser lido com reserva; fica
registrado para a próxima iteração o uso de uma família com massa
perto de zero para esse indicador.

\begin{figure}[H]
\centering
\subfig{fig_est_04_qq_mort_zib.png}{Mortalidade geral, modelo Beta com
inflação em zero}
\subfig{fig_est_04_qq_pct_alta_zoib.png}{Fração de alta complexidade,
modelo ZOIB}
\caption{Resíduos quantílicos dos modelos de taxa.}
\label{fig:qqprimeira}
\end{figure}

\begin{figure}[H]
\centering
\subfig{fig_est_04_qq_producao_bn.png}{Produção, Binomial Negativa}
\subfig{fig_est_04_qq_ocup_uti_pos.png}{Ocupação de UTI entre os
positivos, o único ajuste insatisfatório}
\caption{Resíduos quantílicos da produção e da parte positiva da UTI.}
\label{fig:qqsegunda}
\end{figure}

\comoler{quanto mais os pontos azuis acompanham a linha tracejada,
melhor o modelo descreve os dados; o valor de r no canto é essa
aderência resumida em um número (1 seria perfeito). Compare o
alinhamento quase exato da fração de alta complexidade com o descolamento
da cauda esquerda na UTI, que motivou a ressalva do texto.}

Fecham o quadro as ressalvas permanentes de leitura. A identificação
do efeito de gestão dentro de hospital repousa em \textbf{cinco
conversores}, três deles com um único ano após a mudança; as margens
de erro agrupadas com tão poucas unidades tratadas tendem a ser
otimistas, e os efeitos devem ser lidos como ordens de grandeza. A
categoria Privado (3 CNES) permaneceu como controle de absorção em
todos os modelos, sem interpretação. E os resultados de mortalidade e
ocupação do estudo de eventos carregam a contaminação de tendências
anteriores já descrita, que o tempo de permanência não apresenta.
